\chapter{Evaluación sobre datos reales}
\label{chap:datos_reales}

El segundo bloque de resultados utiliza ECG clínicos reales cedidos por el Hospital Universitario Central de Asturias (HUCA) dentro del recurso Brugada HUCA \cite{morislara2026brugada}. La anotación procede de la recogida paciente a paciente y de la revisión manual especializada. La finalidad es medir cómo se comporta la CNN cuando abandona el simulador y se evalúa sobre trazados clínicos reales. Este paso es necesario porque un buen resultado sintético no garantiza utilidad en un dominio con ruido de adquisición, variabilidad clínica y una etiqueta definida a nivel de paciente.

\section{Cohorte final}

El procesado parte de 363 pacientes. Se excluyen 46 casos con patrón basal limítrofe que no deben mezclarse como positivos o negativos en la tarea principal. La cohorte final contiene 317 pacientes: 116 casos Brugada confirmados y 201 controles. Para cada paciente se generan tres imágenes, una por V1, V2 y V3, lo que produce 951 imágenes.

\begin{table}[H]
  \centering
  \caption{Cohorte cedida por el HUCA utilizada en la evaluación CNN.}
  \label{tab:huca_cohorte_final}
  \begin{tabular}{lrr}
    \toprule
    Grupo & Pacientes & Imágenes \\
    \midrule
    Brugada confirmado & 116 & 348 \\
    Control & 201 & 603 \\
    Excluido por patrón basal limítrofe & 46 & -- \\
    \midrule
    Total evaluado & 317 & 951 \\
    \bottomrule
  \end{tabular}
\end{table}

La etiqueta real disponible no es equivalente a las tres etiquetas morfológicas del simulador. Los datos cedidos por el HUCA no contienen anotaciones independientes de RBBB, elevación ST e inversión T para cada imagen. Por ello, el experimento real se formula como caso Brugada confirmado frente a control y sus conclusiones se limitan a triaje de esa variable clínica. Es importante recordar que en la práctica asistencial, un resultado positivo del modelo no equivale a un diagnóstico: señala que el trazado debería ser revisado por un especialista y, si procede, iniciar el circuito diagnóstico que puede incluir prueba farmacológica con bloqueantes del canal de sodio \cite{wilde2023scbt}.

\section{Preprocesamiento}

La preparación selecciona las derivaciones precordiales derechas, detecta un latido central y recorta una ventana de 300 ms antes de R y 600 ms después de R. La salida conserva el identificador de paciente, derivación, etiqueta clínica, ruta de imagen y plan de pliegues LOOCV. El objetivo es obtener una representación visual comparable a la sintética sin borrar la variabilidad propia del ECG real.

\begin{figure}[H]
  \centering
  \safeincludegraphics[width=0.96\textwidth]{results/fig_huca_preprocesado.pdf}
  \caption{Preprocesamiento de los ECG cedidos por el HUCA para V1, V2 y V3.}
  \label{fig:huca_preprocesado}
\end{figure}

La separación por paciente es obligatoria. Si V1, V2 y V3 de un mismo paciente se repartieran entre entrenamiento y test, el modelo podría aprovechar características de adquisición o escala en lugar de aprender rasgos clínicamente relevantes.

\section{Resultados CNN base}

La tabla \ref{tab:huca_baseline_final} resume la CNN base entrenada y evaluada sobre la cohorte cedida por el HUCA. A diferencia del dominio sintético, aumentar \(K\) no produce una mejora monotónica. El mejor punto de exactitud equilibrada aparece con \(K=16\), pero solo alcanza 0,516. Con \(K=32\) la exactitud equilibrada baja a 0,499. Estos valores son esencialmente indistinguibles del azar.

\begin{table}[H]
  \centering
  \caption{Resultados CNN ResNet18 sobre los datos cedidos por el HUCA.}
  \label{tab:huca_baseline_final}
  \small
  \begin{tabular}{rrrrrr}
    \toprule
    \(K\) & BA & F1 & Sens. & Esp. & ROC AUC \\
    \midrule
    2  & \(0,509 \pm 0,010\) & 0,484 & 0,693 & 0,325 & 0,525 \\
    4  & \(0,509 \pm 0,023\) & 0,493 & 0,730 & 0,289 & 0,495 \\
    8  & \(0,490 \pm 0,018\) & 0,469 & 0,672 & 0,308 & 0,496 \\
    16 & \(0,516 \pm 0,023\) & 0,480 & 0,658 & 0,375 & 0,569 \\
    32 & \(0,499 \pm 0,012\) & 0,466 & 0,644 & 0,355 & 0,491 \\
    \bottomrule
  \end{tabular}
\end{table}

\begin{figure}[H]
  \centering
  \safeincludegraphics[width=0.80\textwidth]{results/cnn_huca_balanced_accuracy_by_k.png}
  \caption{Exactitud equilibrada de la CNN base sobre los datos cedidos por el HUCA.}
  \label{fig:cnn_huca_ba}
\end{figure}

La sensibilidad es moderada, pero la especificidad es baja. En el mejor punto \(K=16\), la matriz agregada de tres semillas contiene 229 verdaderos positivos, 226 verdaderos negativos, 377 falsos positivos y 119 falsos negativos. El modelo tiende a sobredetectar casos Brugada, lo que podría ser tolerable únicamente como alerta preliminar extremadamente conservadora, pero no como herramienta de descarte ni de diagnóstico.

\begin{figure}[H]
  \centering
  \safeincludegraphics[width=0.58\textwidth]{results/cnn_huca_k16_confusion_matrix.png}
  \caption{Matriz de confusión agregada de la cohorte cedida por el HUCA para \(K=16\).}
  \label{fig:cnn_huca_confusion}
\end{figure}

Este resultado es el hallazgo más importante de la línea CNN: con tan pocos datos, el entrenamiento supervisado de una ResNet18, incluso con pesos preentrenados de ImageNet, no consigue extraer señal útil de ECG clínicos reales de Brugada. La red aprende en el simulador porque las imágenes son limpias y las etiquetas derivan de reglas controladas. En el dominio real, la variabilidad clínica, el ruido de adquisición y la complejidad del fenotipo superan la capacidad del modelo.

\section{Brecha sintético-real}

La comparación directa muestra una brecha de dominio marcada. En el simulador la exactitud equilibrada aumenta hasta 0,848 con \(K=32\). En los datos cedidos por el HUCA, el modelo se mantiene alrededor de 0,50. Esto indica que la señal visual aprendida en condiciones controladas no se transfiere de forma directa a la variabilidad clínica.

\begin{figure}[H]
  \centering
  \safeincludegraphics[width=0.86\textwidth]{results/cnn_balanced_accuracy_sim_vs_real.png}
  \caption{Comparación de exactitud equilibrada entre el conjunto sintético y la cohorte cedida por el HUCA.}
  \label{fig:sim_vs_real_ba}
\end{figure}

\begin{figure}[H]
  \centering
  \safeincludegraphics[width=0.86\textwidth]{results/cnn_f1_sim_vs_real.png}
  \caption{Comparación de F1 entre el conjunto sintético y la cohorte cedida por el HUCA.}
  \label{fig:sim_vs_real_f1}
\end{figure}

La brecha tiene varias causas plausibles: ruido y variabilidad de adquisición, diferencias de escala, presencia de múltiples fenotipos, etiqueta clínica no reducible a una sola ventana de ECG y ausencia de anotaciones morfológicas finas en la cohorte real. El resultado no invalida el simulador; delimita su alcance y confirma que el problema de pocos datos no se resuelve con un simulador más limpio.

\section{Inspección cualitativa}

El Grad CAM real confirma que las activaciones pueden localizarse sobre regiones del trazado, pero también muestra la dificultad de interpretar modelos entrenados con pocos ejemplos sobre datos clínicos heterogéneos. La atención visual no corrige por sí sola la falta de especificidad ni resuelve la diferencia entre etiqueta clínica y morfología observable.

\begin{figure}[H]
  \centering
  \safeincludegraphics[width=0.92\textwidth]{results/cnn_huca_gradcam_example.png}
  \caption{Ejemplo Grad CAM del modelo CNN sobre los datos cedidos por el HUCA.}
  \label{fig:cnn_huca_gradcam}
\end{figure}

\section{Lectura de la evaluación real}

La evaluación sobre los datos cedidos por el HUCA delimita el alcance de la línea CNN. La red convolucional, incluso con pesos preentrenados, no logra un rendimiento fiable en datos reales de Brugada con tan pocos ejemplos etiquetados. La exactitud equilibrada se mantiene cerca del azar y la especificidad es demasiado baja para cualquier uso práctico más allá de investigación exploratoria.

Este resultado marca la referencia que debe superar o matizar la línea VLM/ICL. La siguiente fase debe determinar si un modelo de visión y lenguaje, con conocimiento visual y médico adquirido durante el preentrenamiento, ofrece una mejor relación entre rendimiento, coste y cantidad de datos etiquetados cuando recibe esos mismos \(K\) ejemplos como demostraciones en contexto en lugar de como conjunto de entrenamiento.
